% Generated draft: requires mathematical and source review.
\documentclass[11pt]{article}
\usepackage[margin=25mm]{geometry}
\usepackage{fontspec}
\setmainfont{Latin Modern Roman}
\usepackage{amsmath,amssymb,mathtools,unicode-math}
\setmathfont{Latin Modern Math}
\usepackage{enumitem,needspace,xurl,hyperref}
\setlength{\parindent}{0pt}
\setlength{\parskip}{5pt}
\setlength{\emergencystretch}{3em}
\newcommand{\N}{\mathbb{N}}
\newcommand{\Z}{\mathbb{Z}}
\newcommand{\Q}{\mathbb{Q}}
\newcommand{\R}{\mathbb{R}}
\newcommand{\C}{\mathbb{C}}
\newcommand{\F}{\mathbb{F}}
\newcommand{\A}{\mathbb{A}}
\newcommand{\PP}{\mathbb{P}}
\newcommand{\E}{\mathbb{E}}
\newcommand{\eps}{\varepsilon}
\newcommand{\dd}{\,\mathrm{d}}
\begin{document}
\section*{UnsolvedMath ID: 10600050}
\textbf{Source problem code:} catalogue entry
\subsection{Definitions and mathematical statement}
A free-knot diagram is conventionally a framed \(4\)-valent graph with one unicursal component, where the framing specifies opposite half-edges at each vertex; diagrams are taken modulo free Reidemeister moves. Let \(D\) represent the given free knot and let \(W\) be its given cobordism, in the source's unspecified cobordism category.

A usual \(\Z/2\Z\)-valued parity assigns each vertex \(v\) a label \(p(v)\in\Z/2\Z\), compatibly with moves: a vertex removed by a first move has label zero, the two vertices of a second move have equal labels, and corresponding labels in a third move agree with the sum of the three involved labels zero. This explains a standard parity notion; the source does not state its coefficient group or axioms.

Problem-list item 50 asks to construct a parity on the given free knot using only the given cobordism and respecting only moves inside that cobordism. Schematically, it seeks labels \(p_W(v)\) intrinsic to the specified cobordism data and compatible with the induced, permitted changes of sections/diagrams in \(W\). It does not require compatibility with every unrelated cobordism or every unrestricted boundary move. The permitted interior moves, the section structure, and the equivalence on \(W\) are not defined in this extract; even whether \(W\) is a null-cobordism or has another boundary knot is unspecified.

The record gives no nontriviality or detection requirement. Since an identically zero assignment satisfies the usual parity axioms, a precise intended research problem needs additional source criteria distinguishing the sought cobordism-dependent construction. No such criterion or particular geometric rule is invented here.
\subsection{Short English statement}
Use a specified free-knot cobordism itself to define parity labels, compatible with the moves allowed inside it. The cobordism and section-move conventions are missing, as is a condition excluding a vacuous parity. The supplied assessment describes an ongoing parity-cobordism program without documenting this precise construction as completed.
\subsection{Sources}
\begin{itemize}
\item Local UnsolvedMath registry, permanent ID \(10600050\), exact source code \texttt{catalogue entry}, virtual-knot problem-list item 50, ``The problem about cobordisms in sections.'' No individual proposer or proposed year is supplied; the title is truncated.
\item \emph{Unsolved Problems in Virtual Knot Theory and Combinatorial Knot Theory}, revised survey, dedicated problem-list item 50: \url{https://arxiv.org/abs/1409.2823}. Background: 60 top-level items, compound items retained intact; source-TeX extraction, accessed 2026-07-29.
\item Supplied incomplete references: ``Parity and Cobordisms of Free Knots'' (2010), arXiv:1001.xxxx; ``Cobordisms of Free Knots and Gauss Words'' (2009), arXiv:0904.xxxx; ``A Sliceness Criterion for Odd Free Knots'' (2017), arXiv:1707.xxxx. The xxxx identifiers are placeholders and have not been completed.
\item Source evidence NEEDS\_REVIEW presents an open problem and calls for a current primary-literature audit. Supplied OPEN-TRIAGE/borderline summary says the parity-cobordism program exists and may address the construction, but completion of this particular request is not documented as a closed problem.
\item Dataset attribution: Ulam AI / UnsolvedMath Contributors, CC BY 4.0. Source rights NEEDS\_REVIEW: publicly accessible author/publisher copy, redistribution terms requiring review.
\end{itemize}
\end{document}
